\documentclass[a4paper,12pt]{article}

\usepackage[serbian]{babel}
\usepackage[utf8]{inputenc}
\usepackage[T1]{fontenc}
\usepackage{lmodern}
\usepackage{microtype}
\usepackage[a4paper,margin=2.5cm]{geometry}
\usepackage{setspace}
\usepackage{parskip}
\onehalfspacing
\setlength{\parindent}{0pt}

\usepackage{graphicx}
\usepackage{float}
\usepackage{booktabs}
\usepackage{hyperref}
\usepackage{cite}
\usepackage{xcolor}

\usepackage{fancyhdr}
\setlength{\headheight}{15pt}
\pagestyle{fancy}
\fancyhf{}
\fancyhead[L]{\small \nouppercase{\leftmark}}
\fancyfoot[C]{\thepage}

\title{%
\sffamily
\Large Univerzitet u Beogradu\\
\large Matematički fakultet\\[0.8em]
\textbf{Računarstvo i društvo}\\[1.2em]
\rule{\linewidth}{0.5pt}\\[0.6em]
\LARGE \textbf{Kako video igre utiču na malu decu}\\
\large Seminarski rad\\[0.6em]
\rule{\linewidth}{0.5pt}\\[1.2em]
}

\author{%
\sffamily
\begin{tabular}{@{}ll@{}}
\textbf{Student:} & Danilo Nikolaš (112/2021) \\
\textbf{Profesor:} & Dr. Sana Stojanović Durdević \\
\end{tabular}
}

\date{\sffamily Novembar 2025}

\begin{document}
\maketitle
\newpage
\tableofcontents
\newpage


\section{Uvod}

Digitalne tehnologije postale su prirodni deo odrastanja savremene dece. 
Dok starije generacije pamte odrastanje uz dvorišne igre, televiziju i društvene igre, 
današnje detinjstvo se u velikoj meri odvija i u virtuelnim prostorima. 
Video igre, mobilne aplikacije i online platforme spadaju među najčešće korišćene oblike zabave kod dece mlađeg uzrasta.

Međunarodne studije pokazuju da deca između 5 i 10 godina danas provode 
prosečno između 40 i 90 minuta dnevno u igranju video igara, 
pri čemu najveći deo vremena otpada na mobilne igre i jednostavne aplikacije 
dizajnirane za kratke i česte sesije igranja \cite{eu_kids_online_2023}. 
Ovaj fenomen postavlja opravdano pitanje: \textbf{šta zapravo video igre rade razvijajućem mozgu i ponašanju dece?}

Tema je kompleksna — igre nisu same po sebi „dobre” ili „loše”. 
Njihov uticaj zavisi od mnogo faktora: vrste igre, dužine igranja, uzrasta deteta, 
roditeljskog nadzora i socio-emocionalnog okruženja. 
Cilj ovog rada je da pruži uravnotežen pogled zasnovan na dostupnoj literaturi, 
da objasni psihološke i neurološke procese uključene u igranje, te predstavi koristi, rizike i smernice za zdravu upotrebu digitalne zabave.


\section{Razvoj mozga kod dece i interakcija sa digitalnim sadržajima}

Deca u ranom školskom uzrastu intenzivno razvijaju pažnju, sposobnost planiranja, kontrolu emocija i razumevanje pravila i posledica. Ove sposobnosti se ne formiraju odjednom, već postepeno, kroz iskustva u realnom svetu i kroz interakciju sa ljudima.

Video igre su danas deo svakodnevice dece i mogu biti jedan od faktora koji utiču na njihov razvoj. One nude brzo smenjivanje događaja, vizuelno bogate scene i nagrade za brze reakcije. Kada su igre kvalitetne i birane u skladu sa uzrastom deteta, deca mogu razvijati pažnju, logičko razmišljanje, strateško planiranje i koordinaciju pokreta \cite{granic2014benefits}. 

Na primer, igre koje zahtevaju rešavanje zadataka, planiranje poteza ili izgradnju struktura (poput edukativnih i kreativnih naslova) mogu doprineti boljem snalaženju u prostoru, bržoj obradi informacija i razvoju kreativnosti. Takođe, narativne igre ili igre u kojima se sarađuje sa drugima mogu pomoći detetu da razume emocije, odnose i saradnju.

Sa druge strane, ako dete previše vremena provodi u brzim i veoma stimulativnim igrama, može imati poteškoća da se fokusira na sporije aktivnosti koje zahtevaju strpljenje, poput čitanja, pisanja domaćeg zadatka ili dužeg slušanja u školi. Istraživanja ukazuju da prekomerna izloženost brzom digitalnom sadržaju može smanjiti toleranciju na dosadu i otežati održavanje pažnje \cite{anderson2010violent}.


\subsection*{Aktivna i pasivna upotreba tehnologije}
Važno je razlikovati aktivnu od pasivne upotrebe digitalnih uređaja. 
Aktivna upotreba je kada dete rešava zadatke, kreira sadržaj, istražuje, planira i komunicira, ima pozitivniji uticaj na razvoj. 
Suprotno tome, pasivno gledanje sadržaja ili nasumično „kliktanje“ bez jasnog cilja uglavnom ne doprinosi učenju, a može dovesti do površne pažnje i brže frustracije.

\subsection*{Jezički i socijalni razvoj}
Igre mogu biti izvor novih reči i pojmova, naročito na engleskom jeziku, što pomaže razvoju vokabulara. Međutim, digitalna komunikacija nikako ne može zameniti realnu interakciju sa drugom decom i odraslima, koja je ključna za emocionalni i socijalni razvoj.

\subsection*{Emocionalne reakcije}
Igranje može izazvati snažne emocije — uzbuđenje, ponos, ali i ljutnju ili frustraciju. 
Ako roditelji prate i razgovaraju sa detetom o emocijama tokom igranja, dete može naučiti kako da bolje upravlja svojim reakcijama. 
Bez vođstva, međutim, frustracija i impulsivno ponašanje mogu postati češći.

\bigskip

Na kraju, presudni faktori u oblikovanju uticaja video igara na razvoj deteta jesu umerenost, kvalitet sadržaja i aktivno prisustvo roditelja u digitalnim navikama deteta.



\section{Vrste video igara i mehanizmi nagrađivanja}

Video igre namenjene deci mlađeg uzrasta mogu se podeliti prema strukturi, cilju i nivou interaktivnosti. Različite vrste igara stimulišu različite kognitivne i emocionalne procese, pa izbor igre ima veliki značaj za razvoj deteta.

Najčešće kategorije uključuju:
\begin{itemize}
    \item \textbf{Edukativne igre} 
    (npr. \textit{Kahoot Kids}, \textit{Duolingo}) - podstiču logičko razmišljanje, učenje jezika, matematike i rešavanje problema,
    \item \textbf{Kreativne i sandbox igre}
    (npr. \textit{Minecraft}, \textit{Roblox}) - razvijaju maštu, planiranje, saradnju i sposobnost samostalnog stvaranja sadržaja
    \item \textbf{Akcione i refleksne igre} 
    (npr. \textit{Crash Bandicoot}, \textit{Mario Kart}) - unapređuju brzinu reakcije, koordinaciju oko–ruka i vizuelnu percepciju.
    \item \textbf{Mobilne „casual” igre} 
    (npr. \textit{Cut the Rope}, \textit{Subway Surfers}) - karakterišu kratke sesije i česte nagrade, što može biti zabavno, ali i podsticati naviku čestog vraćanja u igru.
    \item \textbf{Kooperativne online igre} 
    (npr. \textit{Minecraft Multiplayer}, \textit{Brawl Stars}) - omogućavaju zajedničku igru i komunikaciju sa vršnjacima, ali nose i rizik od pritiska grupe i neprimerenog online okruženja bez nadzora.
\end{itemize}

Edukativne igre često koriste princip \textbf{postepenog učenja i nagrađivanja} — dete rešava problem, dobija povratnu informaciju i napreduje. Ovakav pristup podstiče radoznalost, upornost i logičko mišljenje. Kreativne igre omogućavaju deci da stvaraju i primenjuju ideje u virtuelnom okruženju, što jača samostalnost i planiranje.

\subsection*{Instant nagrađivanje i dopaminski ciklusi}
Posebno je važno razumeti mehanizam \textbf{instant nagrađivanja}. Kratke nagrade u vidu bodova, vizuelnih efekata, nivoa i pohvala aktiviraju dopaminski sistem i podstiču ponavljanje ponašanja \cite{przybylski2017motivation}. Kod dece, čije su samokontrola i impulsna regulacija u razvoju, ovaj efekat može biti snažniji.

Zbog toga igre često uvode elemente:
\begin{itemize}
    \item skupljanja nagrada,
    \item dnevnih zadataka,
    \item poena za napredovanje,
    \item ograničenog vremena ili „bonus perioda”.
\end{itemize}

Ovi sistemi mogu doprineti ustrajnosti i motivaciji, ali u mobilnim igrama ponekad vode ka navikavanju deteta na stalnu stimulaciju i potrebu za čestim povratkom u igru.

\subsection*{Unutrašnja naspram spoljašnje motivacije}
Važno je razlikovati:
\begin{itemize}
    \item \textbf{unutrašnju motivaciju} — dete igra jer voli izazov, kreativnost i učenje;
    \item \textbf{spoljašnju motivaciju} — dete igra zbog nagrada, medalja i vizuelnih stimulansa.
\end{itemize}

Igre koje podstiču unutrašnju motivaciju, kao što su kreativne i logičke igre, pozitivno utiču na razvoj. Nasuprot tome, igre sa stalnim nagradama i brzim vizuelnim podsticajima mogu otežati detetu da razvije strpljenje i sposobnost dugoročnog planiranja.

\subsection*{Balans i izbor sadržaja}
Nije svaka igra jednako korisna. Razlika između obrazovne igre i nasumičnog „tapkanja“ na ekranu može biti ogromna. Izbor sadržaja i roditeljsko vođstvo igraju ključnu ulogu u tome da li igra postaje izvor kreativnosti ili izvor frustracije i zavisnosti.

Kombinacija edukativnih, kreativnih i socijalno balansiranih igara, uz ograničeno vreme i nadzor, predstavlja najzdraviji pristup digitalnoj zabavi u ranom uzrastu.


\section{Pozitivni efekti igranja kod dece}

Iako su često predmet zabrinutosti, video igre mogu imati značajne razvojne koristi kada se koriste umereno i uz odgovarajući izbor sadržaja. Brojna istraživanja ukazuju da pravilno odabrane igre doprinose razvoju kognitivnih, socijalnih i emocionalnih veština \cite{granic2014benefits}.

\subsection*{Razvoj kognitivnih sposobnosti}
Igre koje zahtevaju planiranje, rešavanje problema i strateško razmišljanje podstiču aktivno učenje. 
Deca koja povremeno igraju igre koje uključuju logiku, prostornu orijentaciju ili zadatke brzog odlučivanja mogu pokazati bolje rezultate u kognitivnim testovima, bržu obradu informacija i bolju sposobnost rasuđivanja.

Vizuelno–prostorne veštine se takođe značajno razvijaju — igrači razvijaju sposobnost da u glavi zamišljaju predmete iz različitih uglova i lakše se orijentišu u prostoru.

\subsection*{Koordinacija i motoričke veštine}
Adekvatan izbor igara može doprineti razvoju koordinacije oko–ruka, preciznog upravljanja pokretom i finih motornih veština. 
Ovi efekti su posebno izraženi kod igara koje uključuju zadatke poput preciznog ciljanja, upravljanja objektima ili veštine ritma i koordinacije.

\subsection*{Kreativnost i maštovito mišljenje}
Kreativne i „sandbox” igre, poput Minecraft-a, omogućavaju deci da grade, eksperimentišu i razvijaju originalne ideje u virtuelnom okruženju. 
Deca često prenose inspiraciju iz igara u realni svet, kroz crtanje, izradu maketa ili simboličku igru, što podržava razvoj imaginacije i divergentnog mišljenja.

\subsection*{Socijalna interakcija i timski rad}
Uprkos čestoj pretpostavci da igre izoluju decu, određene platforme ohrabruju timski rad, zajedničko planiranje i komunikaciju sa vršnjacima. 
Kooperativne igre mogu pomoći u razvoju veština saradnje, pregovaranja i zajedničkog rešavanja problema. Kod introvertne dece, ovakve interakcije ponekad olakšavaju socijalnu integraciju i izražavanje u grupi.

\subsection*{Emocionalna otpornost i samopouzdanje}
Igre često nude bezbedno okruženje za suočavanje sa neuspehom i učenjem iz grešaka. 
Postepeno prevazilaženje izazova i napredovanje kroz igru može ojačati samopouzdanje i ustrajnost. 
Neka istraživanja ukazuju da deca koja imaju pozitivno iskustvo sa rešavanjem zadataka u digitalnom okruženju pokazuju veći nivo motivacije za rešavanje problema i u realnom svetu.

\subsection*{Učenje kroz simulaciju}
Simulacione i edukativne igre omogućavaju deci da istražuju koncepte iz nauke, matematike, jezika ili društvenih nauka na interaktivan način. 
Kada su pažljivo odabrane, ovakve igre mogu nadograditi učenje iz škole, razviti radoznalost i ponuditi bezbedno okruženje za eksperimentisanje i donošenje odluka.

\bigskip

Kada su dobro odabrane i umereno korišćene, igre mogu biti koristan alat u razvoju deteta.


\section{Rizici i negativni efekti}

Iako video igre mogu imati značajne razvojne koristi, nepravilna ili preterana upotreba povezana je sa određenim rizicima. Efekti najčešće zavise od dužine igranja, vrste sadržaja, temperamenta deteta i prisustva roditeljske kontrole.

\subsection*{Pažnja i koncentracija}
Dugotrajna izloženost brzim vizuelnim stimulacijama može smanjiti toleranciju deteta na sporije i manje dinamične aktivnosti, poput učenja ili čitanja. 
Prema pojedinim studijama, deca koja provode mnogo vremena u brzim akcijskim igrama mogu imati poteškoće u zadržavanju pažnje i prelasku na zadatke koji zahtevaju dužu koncentraciju \cite{anderson2010violent}. 
Ovaj efekat je naročito izražen kod dece sa već postojećim teškoćama u samokontroli.

\subsection*{Emocionalna regulacija i frustracija}
Brza smena uspeha i neuspeha u igrama može izazvati snažne emocionalne reakcije. 
Mlađa deca često nemaju dovoljno razvijene mehanizme za kontrolu emocija, pa se frustracija može ispoljiti kroz plač, ljutnju ili impulsivne reakcije prilikom prekida igre ili gubitka.
Ukorenjivanje ovakvih reakcija može otežati razvoj emocionalne stabilnosti i tolerancije na stres.

\subsection*{San i dnevni ritam}
Igranje video igara neposredno pre spavanja može dovesti do teškoća sa uspavljivanjem zbog povećane mentalne pobuđenosti i izlaganja plavom svetlu. 
Prema istraživanjima, deca koja imaju stalni pristup ekranima noću pokazuju kraći i fragmentisan san, kao i povećanu razdražljivost tokom dana \cite{bold2022sleep}. 
Narušen san utiče na koncentraciju, raspoloženje i učenje.

\subsection*{Fizička aktivnost i zdravlje}
Prekomerno vreme provedeno u sedenju uz ekran može zameniti fizičku aktivnost, što doprinosi pasivnijem načinu života i povećanom riziku od gojaznosti, kao i nekih drugih bolesti. 
Svetska zdravstvena organizacija preporučuje ograničeno vreme pred ekranom i svakodnevnu fizičku aktivnost kao deo zdravog razvoja deteta \cite{who_guidelines_2020}.

\subsection*{Socijalna izolacija i ponašanje}
Iako igre mogu biti društvena aktivnost, u nekim slučajevima deca se povlače iz realnih socijalnih interakcija i osećaju nelagodu u kontaktu licem u lice. 
Pojedini sadržaji koji uključuju agresivno ponašanje mogu doprineti povećanoj razdražljivosti i smanjenoj empatiji kod osetljive dece, naročito bez nadzora i razgovora o sadržaju.

\bigskip

Važno je naglasiti da negativni efekti nisu isključivo posledica samih igara, već načina njihove upotrebe. 
Uz jasna pravila, većina navedenih rizika može biti značajno ublažena ili potpuno izbegnuta.


\section{Digitalna zavisnost kod dece i dopaminski mehanizmi}

Digitalna zavisnost predstavlja stanje u kojem dete oseća snažnu potrebu da koristi igru ili digitalni uređaj, često po cenu drugih aktivnosti, socijalnih odnosa ili odmora. Kod male dece ova pojava je posebno izražena jer se izvršne funkcije i sposobnost kontrole impulsa tek razvijaju, dok dopaminski sistem snažno reaguje na nagrade i stimulaciju.

\subsection*{Dopaminski mehanizmi i nagrađivanje}
Igre često koriste kratke cikluse uspeha, vizuelne i zvučne nagrade, otključavanje nivoa i sistem napredovanja pomoću poena za napredovanje, kojima otključavaju nove nivoe, likove i sposobnosti. Ovakvi mehanizmi aktiviraju dopaminske puteve u mozgu, što stvara osećaj prijatnosti i motiviše dete da se vraća igri \cite{pontes2020insights}.  
Što je nagrada brža i češća, veći je rizik da dete postane fokusirano na trenutno zadovoljstvo umesto na dugoročne ciljeve i aktivnosti koje zahtevaju trud i strpljenje.

Kod dece starosti od 5 do 10 godina, čije su sposobnosti samokontrole još u razvoju, ova vrsta stimulacije može dovesti do obrasca ponašanja u kojem dete traži sve bržu i češću nagradu, što povećava rizik od problematičnog korišćenja uređaja.

\subsection*{Rizični obrasci ponašanja}
Rani znaci razvoja problematičnog odnosa prema igrama uključuju:
\begin{itemize}
    \item izraženu frustraciju ili tugu pri prekidu igranja,
    \item izbegavanje drugih aktivnosti u korist ekrana,
    \item smanjenu sposobnost igre bez digitalnih podsticaja,
    \item nervozu, razdražljivost ili povlačenje kada pristup nije moguć,
    \item narušene navike spavanja i dnevnog ritma.
\end{itemize}

Ovi simptomi ne moraju nužno ukazivati na zavisnost, ali predstavljaju signal za roditelje da uvedu jasniju strukturu i pravila korišćenja.

\subsection*{Razlika između interesovanja i zavisnosti}
Važno je razlikovati normalno interesovanje od zavisničkog ponašanja.  
Dete može pokazivati veliko oduševljenje igrom, ali zavisnost se javlja kada:
\begin{itemize}
    \item igra postaje primarna aktivnost dana,
    \item dete gubi kontrolu nad vremenom provedenim u igri,
    \item negativne emocije nastaju kada se igra prekine,
    \item igre se koriste kao glavni način suočavanja sa stresom ili dosadom.
\end{itemize}

Drugim rečima, problem nastaje kada igranje više nije izbor već osećana potreba.

\subsection*{Prevencija i uloga odraslih}
Najefikasnija prevencija uključuje:
\begin{itemize}
    \item izbor adekvatnog sadržaja,
    \item razgovor sa detetom o emocijama tokom igranja,
    \item podsticanje fizičkih i kreativnih aktivnosti van ekrana,
    \item zajedničko igranje i nadzor,
    \item uvođenje dogovorenih vremenskih pravila, uz fleksibilnost i veću kontrolu kada dete postaje previše zaokupirano igrom.
\end{itemize}


\bigskip

Razumevanje načina na koji digitalni sadržaji utiču na ponašanje i navike deteta pomaže roditeljima i školama da pravovremeno reaguju i podrže dete u razvijanju samokontrole i odgovornih digitalnih navika.


\section{Podaci iz istraživanja}

Razumevanje uticaja video igara na decu zahteva oslanjanje na empirijske podatke, a ne samo na teorijske pretpostavke. 
Brojna međunarodna istraživanja ukazuju da digitalna igra predstavlja uobičajen deo detinjstva, ali i da obrasci korišćenja značajno variraju u zavisnosti od uzrasta, načina nadzora i vrste igara.

Prema izveštaju UNICEF-a iz 2023. godine, većina dece između 6 i 12 godina igra video igre barem povremeno, dok trećina mladih igrača pokazuje emocionalnu nelagodu kada im se ograniči pristup igranju \cite{unicef2023gaming}. 
Ovi podaci ukazuju na široku rasprostranjenost digitalne zabave, ali i na potrebu za pažljivim nadzorom nad njenom ulogom u svakodnevnom životu deteta.

\subsection*{Učestalost igranja}
Podaci pokazuju da učestalost i emocionalna reakcija na ograničenje ekrana rastu sa uzrastom. 
Starija deca češće igraju duže seanse, imaju veći pristup zasebnim uređajima i češće koriste online multiplayer okruženja, što utiče na socijalne navike i formiranje digitalnih zajednica.

\begin{table}[H]
\centering
\begin{tabular}{lcc}
\toprule
Starosna grupa & Procenat dece koja igraju & Emocionalna reakcija pri uskraćivanju \\
\midrule
3--5 godina & 52\% & 18\% \\
6--8 godina & 71\% & 25\% \\
9--12 godina & 84\% & 33\% \\
\bottomrule
\end{tabular}
\caption{Učestalost igranja i emocionalna reakcija kod dece \cite{unicef2023gaming}.}
\end{table}

\subsection*{Trendovi korišćenja}
Istraživanja Evropske komisije ukazuju da mlađa deca najčešće igraju kratke mobilne igre, dok se stariji uzrasti okreću kompleksnijim, narativnim i online kooperativnim igrama \cite{eu_kids_online_2023}. 
Ovaj prelaz od jednostavne zabave ka organizovanom digitalnom druženju predstavlja važan razvojni korak, jer menjaju motivaciju — od individualne zabave ka socijalnoj interakciji.

\subsection*{Tumačenje statističkih podataka}

Statistički podaci pokazuju da su video igre postale uobičajen deo detinjstva. Međutim, sama činjenica da dete igra igre ne predstavlja rizik. Istraživanja ne ukazuju da svaki mali igrač razvija negativne posledice — presudni su način korišćenja, kvalitet sadržaja i uključenost odraslih.

Da bi se rezultati pravilno razumeli, važno je posmatrati širi kontekst, jer broj sati nije jedini ni najpouzdaniji pokazatelj:
\begin{itemize}
    \item reakcije deteta i način na koji koristi igru često govore više od same minutaže,
    \item pozitivni efekti se javljaju kada igre imaju edukativnu ili kreativnu vrednost,
    \item balans sa fizičkom aktivnošću, snom i socijalnom interakcijom ostaje ključan.
\end{itemize}

\bigskip

Drugim rečima, razlika nije u tome da li dete igra, već kako, šta igra i koliko je odraslih uključeno u usmeravanje digitalnog iskustva.



\section{Uloga roditelja, obrazovnog sistema i društveni aspekt}

U digitalnom okruženju, roditelji i obrazovne institucije imaju ključnu ulogu u usmeravanju dečjih navika i formiranju zdravog odnosa prema tehnologiji. 
Sama zabrana ili ograničavanje vremena retko daje dugoročne rezultate — najvažniji faktor je \textbf{aktivna uloga odraslih} u oblikovanju digitalnog iskustva deteta.

\subsection*{Roditeljska uloga i primer}
Deca najčešće uče posmatranjem. Ukoliko roditelji sami provode mnogo vremena uz uređaje, teško je očekivati da dete poštuje ograničenja. 
Pozitivan model uključuje: zajedničko vreme bez ekrana, telefon van stola za ručavanje i uređaje izvan spavaće sobe. Vizuelni rasporedi i dogovor unapred smanjuju konflikte i pomažu detetu da razvije samokontrolu.

\subsection*{Zajedničko korišćenje tehnologije}
Najbolji efekti se postižu kada odrasli ne posmatraju digitalnu igru samo kao pasivni nadzor, već povremeno učestvuju u njoj. 
Zajedničko igranje ili razgovor o sadržaju pomaže detetu da razume motive likova, razvija empatiju i učvršćuje porodičnu komunikaciju. 
Umesto kontrole „odozgo“, stvara se partnerska atmosfera, gde dete uči da prepoznaje sopstveno ponašanje i emocije.

\subsection*{Uloga škole i digitalne pismenosti}
Obrazovne institucije sve češće uvode teme digitalne bezbednosti, razvoja kritičkog mišljenja i balansiranja online aktivnosti. 
Škole imaju odgovornost da podrže zdravo korišćenje tehnologije, ali i da pokažu kako digitalni alati mogu biti edukativni, a ne samo izvor zabave.

\subsection*{Etika i uticaj industrije}
Industrija video igara sve više koristi principe \textit{behavioral design} pristupa, uključujući mehanizme koji podstiču duže zadržavanje korisnika u aplikaciji \cite{zagal2021dark}. 
Dok je to prihvatljivo kod odraslih korisnika, etičko pitanje nastaje kada se ovakvi sistemi primenjuju na decu, koja nemaju razvijene mehanizme samokontrole i kritičkog rasuđivanja.

Roditelji, škole i društvo u celini treba da promovišu odgovorno korišćenje digitalnih medija, dok kreatori sadržaja imaju moralnu obavezu da vode računa o najmlađim korisnicima i razvijaju alate koji podstiču kreativnost, a ne zavisničko ponašanje.

\bigskip

Digitalno okruženje nudi brojne razvojne mogućnosti, ali zahteva \textbf{informisan, aktivan i balansiran pristup odraslih}. 
Cilj nije ograničiti pristup tehnologiji, već pomoći detetu da izgradi zdrave navike, kritički pristup sadržaju i sposobnost odgovorne upotrebe digitalnih alata.


\section{Zaključak}

Video igre danas čine sastavni deo detinjstva i odrastanja, i umesto da ih posmatramo kao pretnju, potrebno je razumeti njihov uticaj i pronaći način da se koriste na zdrav i uravnotežen način. Cilj nije zabrana, već stvaranje uslova u kojima igre mogu biti izvor zabave, učenja i razvoja, a ne problem.

Najvažnije smernice izdvojene ovim radom uključuju:
\begin{itemize}
    \item umerenost i balans između igranja, fizičke aktivnosti, sna i školskih obaveza,
    \item aktivno učešće roditelja u izboru sadržaja i praćenju igranja,
    \item razvoj digitalne pismenosti i kritičkog razmišljanja od ranog uzrasta,
    \item postepeno učenje kontrole vremena i emocija tokom korišćenja tehnologije.
\end{itemize}

Zabrana nije rešenje, deca treba da nauče \textit{kako} da koriste tehnologiju, a ne da je izbegavaju. Razgovor, dogovor i razumevanje mnogo su delotvorniji od kažnjavanja ili potpunog uskraćivanja.

Kada postoji podrška, jasno definisana pravila i otvorena komunikacija, video igre mogu predstavljati koristan deo odrastanja: prostor za učenje, razvoj kreativnosti, socijalizaciju i sticanje veština važnih za savremeno društvo. Uz pravilno vođstvo, deca mogu odrastati u digitalno pismene, odgovorne i kreativne korisnike tehnologije.

\newpage
\begin{thebibliography}{99}

\bibitem{granic2014benefits}
I. Granic, A. Lobel, R. Engels, 
``The Benefits of Playing Video Games,''
\textit{American Psychologist}, vol. 69, no. 1, pp. 66–78, 2014.
\url{https://www.apa.org/pubs/journals/releases/amp-a0034857.pdf}
Pristupljeno: Oct. 31, 2025

\bibitem{anderson2010violent}
C.~A. Anderson, A.~Shibuya, N.~Ihara, E.~L. Swing, B.~J. Bushman, A.~Sakamoto, and M.~Saleem,
``Violent Video Game Effects on Aggression, Empathy, and Prosocial Behavior in Eastern and Western Countries: A Meta-Analytic Review,''
\textit{Psychological Bulletin}, vol.~136, no.~2, pp.~151--173, 2010.
\url{https://pubmed.ncbi.nlm.nih.gov/20192553/}
Pristupljeno: Oct. 31, 2025

\bibitem{pontes2020insights}
H. M. Pontes and M. D. Griffiths,
``Emerging insights on Internet Gaming Disorder: Conceptual and methodological issues,''
\textit{Addictive Behaviors Reports}, vol. 11, pp. 100–115, 2020.
\url{https://pmc.ncbi.nlm.nih.gov/articles/PMC7244902/}
Pristupljeno: Oct. 31, 2025

\bibitem{unicef2023gaming}
UNICEF,
``Video Games and Children: A Guide for Parents,''
2023.
\url{https://www.unicef.org/thailand/stories/video-games-and-children-guide-parents}
Pristupljeno: Oct. 31, 2025

\bibitem{who_guidelines_2020}
World Health Organization,
``Guidelines on Physical Activity, Sedentary Behaviour and Sleep for Children,'' 2020.
\url{https://www.who.int/publications/i/item/9789241550536}
Pristupljeno: Oct. 31, 2025.

\bibitem{przybylski2017motivation}
A.~K. Przybylski and N.~Weinstein,
``A Large-Scale Test of the Goldilocks Hypothesis: Quantifying the Relations Between Digital-Screen Use and the Mental Well-Being of Adolescents,''
\textit{Psychological Science}, vol.~28, no.~2, pp.~204--215, 2017.
\url{https://journals.sagepub.com/doi/10.1177/0956797616678438}
Pristupljeno: Oct. 31, 2025

\bibitem{bold2022sleep}
BOLD Science,
``The Effects of Screen Exposure on Sleep in Children,'' 2022.
\url{https://boldscience.org/the-lasting-effects-of-television-and-video-games-in-a-childs-bedroom/}
Pristupljeno: Oct. 31, 2025

\bibitem{eu_kids_online_2023}
EU Kids Online,
``Children’s Online Experiences in 2023: Key Findings from EU Kids Online,''
2023.
\url{https://www.lse.ac.uk/media-and-communications/research/research-projects/eu-kids-online}
Pristupljeno: Oct. 31, 2025

\bibitem{przybylski2014gaming}
A. K. Przybylski,
``Electronic Gaming and Psychosocial Adjustment,''
\textit{Pediatrics}, vol. 134, no. 3, pp. e716–e722, 2014.
\url{https://scispace.com/pdf/electronic-gaming-and-psychosocial-adjustment-z8wb3d4hbi.pdf}
Accessed: Oct. 31, 2025

\bibitem{zagal2021dark}
J. Zagal,
``Dark Patterns in Game Design: How Mechanics Influence Player Behaviour,''
\textit{Games and Culture}, 2021.
[Online]. Available: \url{https://files.core.ac.uk/download/pdf/301007767.pdf}
Pristupljeno: Oct. 31, 2025

\end{thebibliography}

\end{document}
